%---------------------------------------------------------------
% SLO: slovenski povzetek
% ENG: slovenian abstract
%---------------------------------------------------------------
\selectlanguage{slovene} % Preklopi na slovenski jezik
\addcontentsline{toc}{chapter}{Povzetek}
\chapter*{Povzetek}

\noindent\textbf{Naslov:} \ttitle
\bigskip

Realnočasovna simulacija tkanine je pomembna za sodobne interaktivne 3D-sisteme, kot so igre, navidezna resničnost in digitalna predstavitev oblačil, vendar na mobilnih napravah zaradi omejenih računskih virov ostaja zahteven problem. Nevronski modeli ponujajo hitrejšo alternativo klasičnim fizikalno osnovanim metodam, njihova primernost za mobilno izvajanje pa še ni dovolj raziskana. V magistrskem delu smo na poenotenem eksperimentalnem cevovodu primerjali modele TailorNet, HOOD in ContourCraft z uporabo podatkovne množice AMASS ter referenčnih simulacij v okolju Maya nCloth. Rezultati so pokazali, da model HOOD dosega najboljše razmerje med natančnostjo, fizikalno stabilnostjo in hitrostjo izvajanja, zato smo ga izbrali za mobilno obravnavo. V nadaljevanju smo obravnavali prenos modela HOOD na operacijski sistem Android in ovrednotili več mobilnih izvajalnih poti. Rezultati so pokazali, da pospešene poti niso prehitele dobro optimizirane procesorske osnove, so pa omogočile identifikacijo ključnih omejitev mobilnega izvajanja in vzpostavitev enotnega primerjalnega cevovoda za ovrednotenje takšnih pristopov.

\subsection*{Ključne besede}
\textit{\tkeywords}
\clearemptydoublepage